\begin{trapbox}
	\begin{description}[style=nextline, leftmargin=2.8em, labelsep=0.5em, itemsep=0.35em]
		\item[\textbf{\textcolor{CTrap}{易错点一（弦长分解缺失）}}]
		      错误地认为三角形周长必须通过联立方程求出 $|PQ|$ 后才能计算，忽略 $|PQ|$ 可以直接分解为 $|PF_1|+|QF_1|$。

		\item[\textbf{\textcolor{CTrap}{正确理解（定义配对）：}}]
		      因为 $F_1$ 是椭圆内部点，直线 $PQ$ 过 $F_1$ 且交椭圆于 $P,Q$ 两点，所以 $F_1$ 在线段 $PQ$ 上，从而
		      \[
			      |PQ|=|PF_1|+|QF_1|.
		      \]
		      因此
		      \[
			      C=|PF_2|+|QF_2|+|PF_1|+|QF_1|,
		      \]
		      恰好构成两个椭圆定义式。

		\item[\textbf{\textcolor{CTrap}{错因分析（几何忽视）：}}]
		      只看到了“弦长”，没有注意到焦点 $F_1$ 位于弦 $PQ$ 的内部，从而丢掉了关键的分解关系。
	\end{description}

	\medskip
	\begin{description}[style=nextline, leftmargin=2.8em, labelsep=0.5em, itemsep=0.35em]
		\item[\textbf{\textcolor{CTrap}{易错点二（长短轴混淆）}}]
		      误认为周长与 $b$ 有关，写成 $4b$ 或 $2\sqrt{a^2-b^2}$ 等形式。

		\item[\textbf{\textcolor{CTrap}{正确理解（仅与长半轴有关）：}}]
		      由推导可知
		      \[
			      C=4a,
		      \]
		      仅与长半轴 $a$ 有关，与短半轴 $b$ 无关，也不能直接写成与焦距 $2c$ 有关的形式。

		\item[\textbf{\textcolor{CTrap}{错因分析（定义不清）：}}]
		      椭圆定义中
		      \[
			      |PF_1|+|PF_2|=2a
		      \]
		      里的 $a$ 是长半轴，不是短半轴，也不是焦半距 $c$。
	\end{description}

	\medskip
	\begin{description}[style=nextline, leftmargin=2.8em, labelsep=0.5em, itemsep=0.35em]
		\item[\textbf{\textcolor{CTrap}{易错点三（定值误解）}}]
		      认为弦的倾斜度改变时周长会变化，试图每次都用坐标法重新计算。

		\item[\textbf{\textcolor{CTrap}{正确理解（恒为定值）：}}]
		      无论弦的斜率如何，只要弦过焦点，三边长度和恒为
		      \[
			      4a,
		      \]
		      与弦的位置和倾斜程度无关。

		\item[\textbf{\textcolor{CTrap}{错因分析（未用定义）：}}]
		      未从椭圆定义整体考量，陷入复杂的联立和弦长计算。
	\end{description}

	\medskip
	\begin{description}[style=nextline, leftmargin=2.8em, labelsep=0.5em, itemsep=0.35em]
		\item[\textbf{\textcolor{CTrap}{易错点四（忽略退化三角形）}}]
		      直接说“任意焦点弦都与另一焦点构成三角形”，没有注意到当 $PQ$ 恰为长轴弦时，$P,F_2,Q$ 三点共线。

		\item[\textbf{\textcolor{CTrap}{正确理解（普通三角形与退化情形）：}}]
		      若 $PQ$ 不是长轴弦，则 $P,F_2,Q$ 构成非退化三角形，周长为 $4a$；若 $PQ$ 是长轴弦，则三点共线，不是普通三角形，但按三边长度和
		      \[
			      |PF_2|+|QF_2|+|PQ|
		      \]
		      理解时，结果仍为 $4a$。

		\item[\textbf{\textcolor{CTrap}{错因分析（边界情况遗漏）：}}]
		      只关注一般位置的焦点弦，忽略了长轴弦这一特殊边界情况。
	\end{description}
\end{trapbox}